\chapter{Complexity Summary}\label{app:complexity}

\begin{table}[ht]\centering\small
\resizebox{\linewidth}{!}{%
\begin{tabular}{@{}lcccccc@{}}\toprule
Structure / operation & access $i$ & insert front & insert back & erase front & erase back & insert at $i$\\\midrule
Array list & 1 & $n$ & 1 & $n$ & 1 & $n-i$\\
Dynamic array (vector) & 1 & $n$ & $1^*$ & $n$ & $1^*$ & $n-i$\\
SLL, head only & $i$ & 1 & $n$ & 1 & $n$ & $i$\\
SLL, head and tail & $i$ & 1 & 1 & 1 & $n$ & $i$\\
DLL, head and tail & $\min(i,n-i)$ & 1 & 1 & 1 & 1 & $\min(i,n-i)$\\
Circular array (queue/deque) & 1 & 1 & 1 & 1 & 1 & ---\\
ArrayDeque & 1 & 1 & 1 & 1 & 1 & $1+\min(i,n-i)$\\
Linked deque with sentinel & $\min(i,n-i)$ & 1 & 1 & 1 & 1 & $\min(i,n-i)$\\\bottomrule
\end{tabular}}
\caption{Costs of the list operations ($\Theta$, worst case; $^*$ amortised, if the array
grows). Insertion or removal \emph{after a known node} is $\Theta(1)$ in linked lists
(also \emph{before} it in a DLL).}
\end{table}

\begin{table}[ht]\centering\small
\begin{tabular}{@{}lll@{}}\toprule
ADT & implementation & push / pop / top (or front)\\\midrule
Stack & array (top at \cc{items[size-1]}) & $\Theta(1)$ each\\
Stack & SLL (top = head) & $\Theta(1)$ each\\
Queue & circular array (\cc{front}, \cc{size}) & $\Theta(1)$ each\\
Queue & SLL (front = head, rear = tail) & $\Theta(1)$ each\\
Queue & array with front at index 0 & enqueue $\Theta(1)$, dequeue $\Theta(n)$\\
Queue & two stacks & $\Theta(1)$ amortised, $\Theta(n)$ worst\\
Stack & one queue & push $\Theta(n)$, pop/top $\Theta(1)$\\
Deque & circular array, DLL, ArrayDeque & $\Theta(1)$ at both ends\\\bottomrule
\end{tabular}
\caption{Stacks, queues and deques.}
\end{table}

\begin{table}[ht]\centering\small
\begin{tabular}{@{}lll@{}}\toprule
Task & Time & Extra space\\\midrule
linear search in a list & $\Theta(n)$ & 1\\
binary search in a sorted array & $\Theta(\log n)$ & 1\\
insert into a sorted array list & $\Theta(\log n)$ search $+\ \Theta(n)$ shift & 1\\
reverse a list, in place & $\Theta(n)$ & 1\\
middle / $k$-th from end of an SLL & $\Theta(n)$, one pass & 1\\
merge two sorted lists (relinking) & $\Theta(n+m)$ & 1\\
splice after a known node (with tail of the donor) & $\Theta(1)$ & 1\\
cycle detection (Floyd) & $\Theta(n)$ & 1\\
balanced brackets, postfix evaluation, infix$\to$postfix & $\Theta(n)$ & $\Theta(n)$\\
next greater / daily temperatures / stock span & $\Theta(n)$ total & $\Theta(n)$\\
is a pop order possible? (stack permutation) & $\Theta(n)$ & $\Theta(n)$\\
Josephus: simulation / recurrence & $\Theta(nk)$ / $\Theta(n)$ & $\Theta(n)$ / 1\\
sliding window maximum (monotonic deque) & $\Theta(n)$ & $O(k)$\\
longest window with $\max-\min\le L$ & $\Theta(n)$ & $O(n)$\\
DP over the last $k$ values (Jump Game VI) & $\Theta(n)$ & $\Theta(n)$\\\bottomrule
\end{tabular}
\caption{Summary of the algorithmic tasks of this volume.}
\end{table}

\para{Frequently needed sums}
\[\sum_{i=1}^{n}i=\frac{n(n+1)}{2},\qquad\sum_{i=0}^{k}2^i=2^{k+1}-1,\]
\[\sum_{i=0}^{\lfloor\lg n\rfloor}2^i<2n,\qquad C_n=\frac{1}{n+1}\binom{2n}{n}:\ 1,1,2,5,14,42.\]

\chapter{Pitfall Checklist}\label{app:pitfalls}

\begin{enumerate}[(i)]
\item \textbf{Empty, one, two.} Run every function mentally on 0, 1 and 2 items, and on a
  full fixed-capacity structure.
\item \textbf{Save before you free or overwrite.} Read \cc{cur->next} into a variable
  before \cc{free(cur)} or before reassigning \cc{cur->next}.
\item \textbf{Every pointer that names an end must follow the end.} \cc{head}, \cc{tail},
  \cc{size}, \cc{mid}, \cc{front}: update all of them, in every case.
\item \textbf{A function that may change the head} needs the address of the head
  (\cc{node **}) or must return the new head.
\item \textbf{Shallow versus deep.} Copying a struct copies pointers, not nodes. A struct
  that \emph{contains} its own sentinel must never be copied at all.
\item \textbf{Do not walk until \cc{NULL} when the structure changes under you}; walk a
  counted number of steps. Read \cc{q.size()}/\cc{st.size()} once before a loop that changes
  the structure.
\item \textbf{Condition order.} Test the pointer before dereferencing it:
  \cc{p != NULL \&\& p->item == x}, \cc{fast \&\& fast->next}.
\item \textbf{Bounds.} Array list positions: legal insert positions are $0..\cc{size}$,
  legal erase/access positions $0..\cc{size}-1$. State whether positions are 0- or
  1-based.
\item \textbf{Modulo in C.} \cc{-1 \% n} is $-1$. Step back with \cc{(i - 1 + n) \% n}.
  Integer division truncates toward zero: \cc{-7/2} is $-3$.
\item \textbf{Stack discipline.} Pop the right operand first. A transfer between two
  stacks reverses the order; two transfers restore it.
\item \textbf{Empty pops.} \cc{pop()} returns nothing: read \cc{top()}/\cc{front()}
  first. Never call \cc{top}, \cc{front}, \cc{back} or \cc{pop} on an empty container;
  test \cc{!st.empty()} first.
\item \textbf{Monotonic structures store indices}, not values, whenever an item can
  expire or an answer must be written to its position. Decide \cc{<} versus \cc{<=} from
  the word ``strictly''.
\item \textbf{Overflow.} Sums such as $n(n+1)/2$ and products overflow \cc{int}; midpoints
  use \cc{lo + (hi - lo) / 2}.
\item \textbf{Half-open ranges} $[\mathit{lo},\mathit{hi})$ make binary search and window
  operations uniform.
\item \textbf{Recursion depth} equals the length of the list; a recursive walk over $10^5$
  nodes can overflow the call stack.
\item \textbf{State the complexity} of your solution in a comment, and the extra space.
\end{enumerate}

\chapter{STL Quick Reference}\label{app:stl}

Table~\ref{tab:stl} in Section~\ref{sec:stl} lists the containers. The snippets below show
the idioms that replace the course library; they are not used anywhere else in this
volume (the headers \cc{iostream} and \cc{using namespace std;} are omitted).

\programfile[language=C++]{notes/stl_demo.cpp}
It prints:
\outputfile{stl_demo}

\begin{table}[ht]\centering\small
\begin{tabular}{@{}ll@{}}\toprule
You want to & stack / queue / deque\\\midrule
add & \cc{st.push(x)} / \cc{q.push(x)} / \cc{d.push\_front(x)}, \cc{d.push\_back(x)}\\
look (not remove) & \cc{st.top()} / \cc{q.front()}, \cc{q.back()} / \cc{d.front()}, \cc{d.back()}\\
remove & \cc{st.pop()} / \cc{q.pop()} / \cc{d.pop\_front()}, \cc{d.pop\_back()}\\
remove into \cc{x} & \cc{x = st.top(); st.pop();} / \cc{x = q.front(); q.pop();}\\
check & \cc{empty()}, \cc{size()} on all three\\
empty structure & \cc{top}/\cc{front}/\cc{pop} undefined: test \cc{empty()} first\\\bottomrule
\end{tabular}
\caption{Stack, queue and deque operations at a glance. The classes \cc{Stack},
\cc{Queue}, \cc{Deque} of Chapters~\ref{ch:stacks}--\ref{ch:deques} use the same names.}
\end{table}
